\subsection{Matching upper bound for the common-upper model}
\label{sec:optional-upper}

We now prove the upper bounds on the two ranges on which capped outcomes
matter.  Besides completing the upper-bound side of the common-upper curve,
the proof records explicitly why an outcome on the optional cap cannot be
treated as an ordinary deferred outcome.  Throughout this subsection,
\[
 C=1+c,
 \qquad
 \varphi=\frac{1+\sqrt5}{2},
\]
and $z_*>1$ is the solution of $z_*-3=\log z_*$.  We also set
\[
 r=\frac{2}{z_*-1}.
\]

For $c\in[r,1/\varphi]$, let $m=m(c)\in[0,1/\varphi]$ be determined by
\begin{equation}
 \operatorname{artanh}m-m
 =1-\frac1c+\frac12\log\!\left(1+\frac2c\right),
 \qquad
 u=1+\frac2c-m.
 \label{eq:opt:upper-mixed-parametrization}
\end{equation}
The right-hand side of the first equation is strictly increasing in $c$,
as is the left-hand side in $m$.  Consequently $m(c)$ is increasing and
$u(c)$ is decreasing.  The endpoints are
\[
 (c,m,u)=\left(\frac1\varphi,\frac1\varphi,\varphi+2\right)
 \quad\text{and}\quad
 (c,m,u)=(r,0,z_*).
\]

\begin{theorem}
\label{thm:opt:optional-middle-upper}
For every $u\in[\varphi+2,\infty]$, the formally specified algorithm
\textnormal{\textsc{AdaptiveThreshold}} from
\cref{alg:adaptive-threshold}
satisfies
\[
 \operatorname{ALG}\le
 \begin{cases}
  (1+c)\operatorname{OPT}+O_u(n),
     & \varphi+2\le u\le z_*,
       \quad c\text{ given by \eqref{eq:opt:upper-mixed-parametrization}},\\[1mm]
  (1+r)\operatorname{OPT}+O(n),
     & u\ge z_*.
 \end{cases}
\]
The constant in the last $O(n)$ is absolute and uniform over the extended
interval $u\in[z_*,\infty]$.
In particular these are upper bounds on the deterministic
size-asymptotic competitive ratio.
\end{theorem}

\noindent\emph{Proof.}

At $u=\infty$ this is \cref{thm:upper}, with the same algorithm and the same
four-endpoint bank.  For the finite-cap calculations below, fix
$u<\infty$.

\paragraph{The algorithm.}
Use \textsc{AdaptiveThreshold}$(J,c)$ from
\cref{alg:adaptive-threshold}; on the high interval use $c=r$.
For the analysis, immediately before a test let $x$ be the number of jobs not
yet tested, $N$ the number of earlier strictly positive outcomes, and $D$ the
number of those outcomes deferred to the final tail.  Thus
\begin{equation}
 y=\frac{D-CN}{x}\le0
 \label{eq:opt:state-y}
\end{equation}
and the threshold is
\begin{equation}
 a_c(y)=
 \begin{cases}
  1, & y\le-1,\\[1mm]
  \displaystyle
  1+\frac12\log\!\left(\frac{c+2}{c-2y}\right),
      & -1\le y\le0.
 \end{cases}
 \label{eq:opt:threshold}
\end{equation}
This is the same threshold as \cref{eq:upper-A}.  On the mixed interval
$c\ge r$, and hence
\[
 a_c(y)\le a_c(0)\le a_r(0)=\frac1r<u-1.
 \label{eq:opt:threshold-below-cap}
\]
The same inequality holds on the high interval, where we use $c=r$.
Thus every outcome in the capped region $[u-1,u]$ is deferred.

The state $y$ never increases.  Indeed, the numerator in
\eqref{eq:opt:state-y} changes by $0$, $-C$, or $-c$ after a zero, a
positive immediate outcome, or a positive deferred outcome, respectively,
while $x$ decreases.  This monotonicity will allow us to allocate every
pair cost when its relevant endpoint becomes known.

\paragraph{Exact pair objective and the five endpoints.}
Write
\[
 \bar p_i=\min\{p_i,u-1\}.
\]
The clairvoyant optional optimum is SPT on jobs of effective lengths
$1+\bar p_i=\min\{1+p_i,u\}$.  For two jobs tested in the order $i<j$,
let
\begin{equation}
 e_{ij}=
 \begin{cases}
  (p_i-p_j)_+,
     &i\text{ is processed immediately},\\
  1+(p_j-p_i)_+,
     &i\text{ is deferred and }j\text{ is immediate},\\
  1,
     &i\text{ and }j\text{ are deferred}.
 \end{cases}
 \label{eq:opt:pair-excess}
\end{equation}
The first line applies whenever $i$ is immediate, regardless of the
action taken on $j$.  Direct completion-time accounting gives
\begin{align}
 \operatorname{ALG}
 &=\sum_i(1+p_i)
   +\sum_{i<j}\bigl(1+\min\{p_i,p_j\}+e_{ij}\bigr),
 \label{eq:opt:alg-pair-decomposition}\\
 \operatorname{OPT}
 &=\sum_i(1+\bar p_i)
   +\sum_{i<j}\bigl(1+\min\{\bar p_i,\bar p_j\}\bigr).
 \label{eq:opt:opt-pair-decomposition}
\end{align}
It is therefore natural to introduce
\begin{equation}
 \Gamma_n=\sum_{i<j}
 \left(e_{ij}+\min\{p_i,p_j\}
       -C\min\{\bar p_i,\bar p_j\}\right).
 \label{eq:opt:pair-objective}
\end{equation}
Equations \eqref{eq:opt:alg-pair-decomposition}--
\eqref{eq:opt:pair-objective} yield the exact identity
\begin{equation}
 \operatorname{ALG}-C\operatorname{OPT}
 =\Gamma_n-c\binom n2
  +\sum_i\bigl[(1+p_i)-C(1+\bar p_i)\bigr].
 \label{eq:opt:gap-decomposition}
\end{equation}
The final sum is nonpositive.  Below the cap its $i$-th summand is
$-c(1+p_i)$; on the cap it is at most $1-cu\le0$.  Thus it is enough
to control $\Gamma_n$.

Fix a symbolic immediate/deferred path.  Coordinatewise convexity in
\eqref{eq:opt:pair-objective} reduces a worst outcome vector to the five
formal endpoint types
\begin{equation}
\begin{array}{c|c|c}
 \text{type}&\text{outcome}&\text{action}\\ \hline
 Z&0&\text{immediate}\\
 E&0^+&\text{immediate}\\
 M&a_c(y)&\text{immediate}\\
 Q&a_c(y)^+&\text{deferred}\\
 U&u&\text{deferred and capped}.
\end{array}
 \label{eq:opt:five-endpoints}
\end{equation}
For an immediate coordinate the two positive endpoints are $0^+$ and
the live threshold; the genuine zero is kept separate because it does not
enter $N$.  A deferred coordinate below $u-1$ moves to the live
threshold or to $u-1$.  On $[u-1,u]$, the offline contribution is flat
and the algorithmic contribution is nondecreasing, so the latter endpoint
moves to $u$.  The symbolic categories, and therefore all later
threshold states, remain unchanged during each coordinate move.  This is
the only endpoint reduction used below.

\paragraph{A local pair allocation.}
Separate the earlier $E$-outcomes from the substantive outcomes
$M,Q,U$.  Let
\[
 P=\#\{\text{earlier }M,Q,U\},\qquad
 E=\#\{\text{earlier }E\},\qquad
 D=\#\{\text{earlier }Q,U\},
\]
and let $K$ be the number of earlier $U$-outcomes.  Among the ordinary
substantive outcomes, let $I,d$ be the numbers of $M,Q$-outcomes and
let $A_I,A_D$ be the sums of their processing values.  Thus
$P=I+d+K$ and $D=d+K$.

At a current ordinary endpoint, with $a=a_c(y)$, allocate the relaxed
local charges
\begin{equation}
 \ell_Z=\ell_E=D,
 \qquad
 \ell_M=D+a(x+D-CP),
 \qquad
 \ell_Q=D+a(D-CP).
 \label{eq:opt:ordinary-local-charges}
\end{equation}
For a capped endpoint the exact charge is instead
\begin{align}
 \ell_U
 &=D+A_D+uK-C\bigl(A_I+A_D+(u-1)K\bigr)
 \notag\\
 &=d-CA_I-cA_D+K h_U,
 \qquad
 h_U=1+u-C(u-1).
 \label{eq:opt:cap-local-charge}
\end{align}
These formulae follow by checking the finitely many ordered endpoint
pairs.  In particular, an earlier immediate outcome already allocated its
online processing delay when it appeared, so a current cap contributes
only $-Cp$ against that outcome.  Against an earlier ordinary deferred
outcome it contributes $p-Cp=-cp$, and against an earlier cap it
contributes $h_U$.  The same check shows
\begin{equation}
 \Gamma_n\le\sum_i\ell_i;
 \label{eq:opt:pair-allocation-dominates}
\end{equation}
the difference is exactly the sum of the nonnegative own-processing terms
allocated by the $M$-endpoints.  Formula
\eqref{eq:opt:cap-local-charge}, rather than a substitution $p=u$ into
$\ell_Q$, is essential before caps have been put into a prefix.

Threshold monotonicity gives $A_I\ge aI$ and $A_D\ge ad$.  Comparing
\eqref{eq:opt:cap-local-charge} with the hypothetical $Q$-charge at the
same state gives the exact identity
\begin{align}
 \ell_U-\ell_Q
 &=C(aI-A_I)+c(ad-A_D)
   +K\,[1-c(u-1-a)]
 \notag\\
 &\le K\,[1-c(u-1-a)].
 \label{eq:opt:cap-versus-q}
\end{align}

\paragraph{The reusable four-endpoint bank.}
The next potential is the parameterized form of the cap-free bank from
\cref{sec:upper}.  The states $Z,E,M,Q$ are the same four ordinary endpoints;
only their target coefficient changes from $r$ to $c$.  At $c=r$ the formulas
below coincide term for term with the obligatory certificate.  Keeping the
parameter visible lets us add the finite-cap reserve without proving a second
ordinary endpoint theorem.
Define
\begin{equation}
 \eta=\frac{D-CP}{x},
 \qquad b=\frac Ex,
 \qquad y=\eta-Cb.
 \label{eq:opt:eta-b}
\end{equation}
For $-1\le y\le0$, put
\begin{align}
 f_c(y)
 &=\frac{c+2}{4}
   +\frac{c-2y}{4}
      \left(\log\frac{c-2y}{c+2}-1\right),
 \label{eq:opt:f-definition}\\
 h_c(y)&=a_c(y)(1+y),
 \notag\\
 G_c(y,b)&=f_c(y)+b h_c(y)+\frac C2b^2.
 \label{eq:opt:G-region-one}
\end{align}
For $y\le-1$, set instead
\begin{equation}
 G_c(y,b)=g_c(\eta),
 \qquad
 g_c(\eta)=\frac{(1+\eta)_+^2}{2C}.
 \label{eq:opt:G-region-two}
\end{equation}
The two definitions are $C^1$ across $y=-1$.  Indeed, feasibility
implies $0\le b\le1/C$ there, and both sides have value $Cb^2/2$ and
derivatives
\[
 (G_y,G_b)=(b,Cb).
\]
They are also nonnegative: $f_c(-1)=0$,
$f_c'=a_c-1\ge0$, and $h_c,b,g_c\ge0$.

Use the base potential
\begin{equation}
 W_{\mathrm{base}}(x,P,E,D)=xD+x^2G_c(y,b).
 \label{eq:opt:base-potential}
\end{equation}
In the raw state $(x,P,E,D,K)$, the endpoint directions are
\begin{equation}
\begin{array}{c|ccccc}
 &x&P&E&D&K\\ \hline
 v_Z&-1&0&0&0&0\\
 v_E&-1&0&1&0&0\\
 v_M&-1&1&0&0&0\\
 v_Q&-1&1&0&1&0\\
 v_U&-1&1&0&1&1.
\end{array}
 \label{eq:opt:state-directions}
\end{equation}
Let $J_T=\ell_T+\nabla W_{\mathrm{base}}\mathbin{\cdot}v_T$, where for
the moment $T\in\{Z,E,M,Q\}$.  In Region I, direct differentiation and
the identities
\begin{align}
 -2f_c+(y-C)f_c'+a_c(1+y)&=0,
 \notag\\
 1-2f_c+(y-c)f_c'+a_cy&=0
 \label{eq:opt:f-hjb-identities}
\end{align}
give
\begin{align}
 \frac{J_Z}{x}
 &=(-2f_c+yf_c')+b(-h_c+yh_c'),
 \notag\\
 \frac{J_E}{x}
 &=b[-h_c+(y-C)h_c'+C],
 \notag\\
 \frac{J_M}{x}
 &=b[-h_c+(y-C)h_c'+a_cC],
 \notag\\
 \frac{J_Q}{x}
 &=b[-h_c+(y-c)h_c'+a_cC].
 \label{eq:opt:region-one-drifts}
\end{align}
All four quantities are nonpositive.  For the first one,
\begin{equation}
 -2f_c+yf_c'=a_c(c-y)-C\le0;
 \label{eq:opt:zero-drift-sign}
\end{equation}
the right-hand side is zero at $y=-1$ and is nonincreasing because
$a_c'(c-y)-a_c\le1-a_c\le0$.  Also $h_c,h_c'\ge0$.  For the remaining
three drifts, observe that
\begin{equation}
 h_c+(c-y)h_c'
 =a_cC+(c-y)(1+y)a_c'\ge a_cC.
 \label{eq:opt:h-sign-identity}
\end{equation}
The $Q$-bracket in \eqref{eq:opt:region-one-drifts} is therefore
nonpositive; the $M$-bracket is the $Q$-bracket minus $h_c'$, and the
$E$-bracket is the $M$-bracket minus $C(a_c-1)$.

In Region II, the same calculation gives
\begin{align}
 J_Z/x=J_E/x&=-2g_c+\eta g_c',
 \notag\\
 J_M/x&=-2g_c+(\eta-C)g_c'+1+\eta,
 \notag\\
 J_Q/x&=1-2g_c+(\eta-c)g_c'+\eta.
 \label{eq:opt:region-two-drifts}
\end{align}
For $-1\le\eta\le0$, substitution of
$g_c'=(1+\eta)/C$ makes these values $-(1+\eta)/C$,
$-(1+\eta)/C$, and $0$, respectively.  For $\eta\le-1$ their signs
are immediate from $g_c=g_c'=0$.  Hence the base bank pays every
ordinary endpoint.

\paragraph{The cap reserve.}
Define
\begin{equation}
 Z(c)=2+\frac1c+\frac12\log\!\left(1+\frac2c\right),
 \qquad
 \delta=Z(c)-u.
 \label{eq:opt:Z-delta}
\end{equation}
On the mixed curve, \eqref{eq:opt:upper-mixed-parametrization} says precisely
that
\begin{equation}
 \delta=\operatorname{artanh}m.
 \label{eq:opt:delta-atanh}
\end{equation}
For $0\le t\le m$, set
\begin{equation}
 \beta(t)=ct\,[\delta-\operatorname{artanh}t],
 \qquad
 \mathcal V(t)=\int_t^m\beta(s)\,ds,
 \label{eq:opt:reserve-integral}
\end{equation}
and extend $\mathcal V$ by zero for $t\ge m$.  Since
$\beta(m)=0$, this extension is $C^1$.  With
\begin{equation}
 q=\frac Kx,
 \qquad
 \mu=\frac{q}{1+q},
 \qquad
 \mathcal H(q)=(1+q)^2\mathcal V(\mu),
 \label{eq:opt:reserve-H}
\end{equation}
define the full potential
\begin{equation}
 W=xD+x^2\bigl(G_c(y,b)+\mathcal H(q)\bigr).
 \label{eq:opt:full-potential}
\end{equation}
It is nonnegative.

For an ordinary endpoint, $K$ is fixed while $x$ decreases, so the
reserve contributes, after division by $x$,
\begin{equation}
 L_0(q)=-2\mathcal H(q)+q\mathcal H'(q).
 \label{eq:opt:L0}
\end{equation}
A cap also increments $K$, and its contribution is
$L_0+\mathcal H'$.  On the active range $\mu\le m$, differentiation of
\eqref{eq:opt:reserve-H} gives
\begin{equation}
 \mathcal H'=2(1+q)\mathcal V-\beta(\mu),
 \qquad
 L_0=-2(1+q)\mathcal V-q\beta(\mu)\le0.
 \label{eq:opt:reserve-derivatives}
\end{equation}
Thus the reserve can only improve every ordinary drift.

For a cap, feasibility $D\le P$ and $P\ge K$ gives
\begin{equation}
 y=\frac{D-C(P+E)}x\le-\frac{cK}{x}=-cq.
 \label{eq:opt:y-cap-bound}
\end{equation}
While $\mu\le m$, we have $q\le m/(1-m)\le\varphi$ and hence
$cq\le1$.  Monotonicity of $a_c$, together with
\eqref{eq:opt:cap-versus-q}, yields
\begin{align}
 \frac{\ell_U-\ell_Q}{x}
 &\le cq\left[\delta-\frac12\log(1+2q)\right]
 =:B(q).
 \label{eq:opt:cap-residual-B}
\end{align}
The change of variables in \eqref{eq:opt:reserve-H} satisfies
\begin{equation}
 \operatorname{artanh}\mu=\frac12\log(1+2q),
 \qquad
 (1+q)\beta(\mu)=B(q).
 \label{eq:opt:B-beta}
\end{equation}
Consequently the cap residual is cancelled exactly:
\begin{equation}
 L_0+\mathcal H'+B(q)=0
 \qquad(\mu\le m).
 \label{eq:opt:reserve-cancellation}
\end{equation}
When $\mu\ge m$, the reserve vanishes.  At
$q_*=m/(1-m)$, the bracket in
\eqref{eq:opt:cap-versus-q}, evaluated at the largest possible threshold
$a_c(-cq_*)$, is zero by
\eqref{eq:opt:delta-atanh} and \eqref{eq:opt:B-beta}.  For $q\ge q_*$,
\eqref{eq:opt:y-cap-bound} only decreases the threshold, so that bracket is
nonpositive.  The fifth Hamiltonian inequality therefore holds for every
state and for arbitrary interleavings of capped and ordinary outcomes.

\paragraph{Initial calibration.}
Initially $y=b=q=0$, and
\begin{equation}
 f_c(0)=\frac12-\frac c4\log\!\left(1+\frac2c\right).
 \label{eq:opt:f-zero}
\end{equation}
Elementary integration in \eqref{eq:opt:reserve-integral} gives
\begin{equation}
 \mathcal V(0)=\frac c2(\delta-m).
 \label{eq:opt:V-zero}
\end{equation}
Indeed,

\[
 \int_0^m t\operatorname{artanh}t\,dt
 =\frac{(m^2-1)\operatorname{artanh}m+m}{2}.
\]
Using \eqref{eq:opt:upper-mixed-parametrization} and
\eqref{eq:opt:delta-atanh}, we obtain the exact calibration
\begin{equation}
 f_c(0)+\mathcal H(0)
 =f_c(0)+\mathcal V(0)=\frac c2.
 \label{eq:opt:initial-calibration}
\end{equation}
Thus $W_{\mathrm{initial}}=cn^2/2$.

\paragraph{Finite jobs, and uniformity on the plateau.}
We spell out the regularity needed to pass from the fluid Hamiltonian to
unit job transitions.  It is important here to restrict attention to
feasible states.  In Region I, $y\ge-1$ implies
\begin{equation}
 -xy=C(P+E)-D=cP+(P-D)+CE\le x.
 \label{eq:opt:feasible-compactness}
\end{equation}
All terms in the middle expression are nonnegative.  Hence
$P/x,D/x,q\le1/c$ and $b\le1/C$, so every normalized derivative of
the Region-I perspective in \eqref{eq:opt:base-potential} is bounded on
the feasible cone.  In Region II the raw base bank is simply
\begin{equation}
 xD+\frac{(x+D-CP)_+^2}{2C},
 \label{eq:opt:region-two-raw-bank}
\end{equation}
which has bounded one-sided Hessians.  The two pieces are $C^1$ at
$y=-1$, as verified above, and the positive-part square is $C^1$ at
$\eta=-1$.

For the reserve $R(x,K)=x^2\mathcal H(K/x)$, on its active cone
$0\le q\le q_*$ one has bounded $\mathcal H''$, and
\begin{equation}
 R_{xx}=2\mathcal H-2q\mathcal H'+q^2\mathcal H'',
 \quad
 R_{xK}=\mathcal H'-q\mathcal H'',
 \quad
 R_{KK}=\mathcal H''.
 \label{eq:opt:reserve-hessian}
\end{equation}
For $q\ge q_*$, $R=0$; moreover
$\mathcal H(q_*)=\mathcal H'(q_*)=0$.  Thus the join is $C^1$ with
bounded one-sided Hessians.  At $x=0$, set $R=0$: for $K>0$ it is
already locally zero, while at $(x,K)=(0,0)$ two-homogeneity gives zero
gradient.  Altogether $W$ is $C^{1,1}$ on the feasible state cone.  On the
mixed interval its regularity constant may depend on the fixed value of $u$.
When $c=r$ and $m=0$, however, the reserve is identically zero and the base
bank depends only on the absolute constant $r$; its regularity constant is
therefore uniform over all $u\ge z_*$.  This includes the last transition
from $x=1$ to $x=0$.

It follows that a unit transition has the Taylor remainder
\begin{equation}
 \ell_i+W_{i+1}-W_i\le O_u(1).
 \label{eq:opt:finite-taylor}
\end{equation}
The endpoint $E=0^+$ is interpreted one-sidedly.  For a fixed symbolic
word, replace all $E$-coordinates by positive values $\epsilon$.  Its
state trajectory is unchanged, and the exact finite pair objective, being
a finite sum of affine and minimum terms, converges as
$\epsilon\downarrow0$.  Equivalently, retain an
$O(\epsilon n^2)$ error and take this limit for each fixed $n$.
Therefore the formal $E$-Hamiltonian transfers to genuine outcome
vectors without an asymptotic loss.

Summing \eqref{eq:opt:finite-taylor}, using $W\ge0$, and applying
\eqref{eq:opt:initial-calibration} gives
\begin{equation}
 \sum_i\ell_i\le\frac c2n^2+O_u(n).
 \label{eq:opt:ell-telescope}
\end{equation}
Together with \eqref{eq:opt:gap-decomposition} and
\eqref{eq:opt:pair-allocation-dominates}, this yields
\begin{align}
 \operatorname{ALG}-C\operatorname{OPT}
 &\le -c\binom n2+\frac c2n^2+O_u(n)
 =O_u(n).
 \label{eq:opt:mixed-upper-finite}
\end{align}

At $u=z_*$, one has $c=r$, $m=0$, and
$Z(r)=z_*$.  The cap reserve is identically zero, while
\eqref{eq:opt:initial-calibration} remains valid.  If $u$ is increased
above $z_*$, the only adverse fifth-endpoint bracket
$1-r(u-1-a_r(y))$ decreases, and the capped diagonal term in
\eqref{eq:opt:gap-decomposition} also decreases.

Here is the promised uniform finite calculation.  The base bank and its
feasible cone contain no $u$, so there is one constant $\kappa$, depending
only on $r$, for which every ordinary or capped transition with $u\ge z_*$
satisfies
\[
 \ell_i+W_{i+1}-W_i\le\kappa.
\]
Because $W_0=rn^2/2$ and $W_n=0$, telescoping gives
\begin{equation}
 \sum_i\ell_i\le\frac r2n^2+\kappa n.
 \label{eq:opt:uniform-ell}
\end{equation}
In the diagonal term of \cref{eq:opt:gap-decomposition}, a value below the
cap contributes $-r(1+p_i)$.  A capped value contributes at most
$1-r u\le0$.  With the convention $\bar p_i=p_i$ at $u=\infty$, the first
formula applies there and the capped endpoint is absent.  Hence the diagonal
sum is nonpositive uniformly on $[z_*,\infty]$.  Combining
\cref{eq:opt:gap-decomposition,eq:opt:pair-allocation-dominates,eq:opt:uniform-ell}
now yields
\begin{align}
 \operatorname{ALG}-(1+r)\operatorname{OPT}
 &\le-r\binom n2+\frac r2n^2+\kappa n
   =\left(\kappa+\frac r2\right)n.
 \label{eq:opt:uniform-high-gap}
\end{align}
Thus, with $B=\kappa+r/2$,
\begin{equation}
 \operatorname{ALG}\le(1+r)\operatorname{OPT}+B n
 \qquad\text{for every }u\in[z_*,\infty],
 \label{eq:opt:uniform-high-upper}
\end{equation}
and $B$ is independent of $u$.  The estimate is genuinely uniform rather
than the result of sending a parameter-dependent remainder to infinity.  This proves
\cref{thm:opt:optional-middle-upper}.  Finally, because $u>1$ on the
present range, every effective offline length is at least one and
$\operatorname{OPT}\ge n(n+1)/2$.  The additive term in
\eqref{eq:opt:mixed-upper-finite} vanishes in the size-asymptotic ratio.

\hfill$\square$

\newcommand{\capbubblingappendix}{%
\section{Independent certificate by bubbling caps to a prefix}
\label{app:cap-bubbling}
For completeness, we give a second proof of the quadratic estimate.  This
argument is also a useful audit of the fifth-endpoint calculation: it never
inserts $p=u$ into an ordinary deferred formula before the caps have
actually been moved.

Work first in normalized fluid mass.  Immediately before an adjacent pair
$XU$, let $E_0$ be the preceding $E$-mass, $P$ the preceding
substantive mass, $D$ its deferred submass, $x$ the remaining mass, and
$w$ the mass of $U$-endpoints strictly after the pair.  Put
\begin{equation}
 y=\frac{D-C(P+E_0)}x,
 \qquad
 \eta=\frac{D-CP}{x}=y+C\frac{E_0}{x}\le0.
 \label{eq:opt:bubble-state}
\end{equation}
Direct differentiation of the exact pair objective
\eqref{eq:opt:pair-objective}, including the mutual pair and all future
caps, gives the two nontrivial commutators
\begin{align}
 UM-MU
 &=1+\left(1+\eta-\frac{Cw}{x}\right)
       (y-c)a_c'(y),
 \label{eq:opt:cap-M-commutator}\\
 UQ-QU
 &=\left(\eta-\frac{cw}{x}\right)
       (y-c)a_c'(y).
 \label{eq:opt:cap-Q-commutator}
\end{align}
Here $UX-XU$ means ``new objective after moving $U$ left'' minus the
old objective.  The second quantity is nonnegative.  For the first, if
$B=1+\eta-Cw/x\le0$, its second term is nonnegative.  If $B>0$, then
$B\le1$, and on $-1<y\le0$
\[
 \frac{(c-y)B}{c-2y}\le1.
\]
Thus \eqref{eq:opt:cap-M-commutator} is nonnegative as well.  Below
$y=-1$, the threshold is flat.  Direct evaluation also gives
\begin{equation}
 UZ-ZU=UE-EU=1.
 \label{eq:opt:cap-zero-commutators}
\end{equation}
The suffix terms involving $w$ are crucial: they make every adjacent
exchange valid even while other caps remain to the right.  Repeated swaps
therefore put all caps into one prefix.

After that prefix there is no future cap.  Let $e,b,d$ be the preceding
$E$-mass, substantive mass, and deferred substantive mass, and define
\[
 y=\frac{d-C(e+b)}x,
 \quad v=\frac ex,
 \quad h=d-Cb,
 \quad k=x+d-Cb.
\]
The four ordinary commutators are
\begin{align}
 ME-EM&=a_c+a_c'(C-y)\frac{k}{x},
 &QE-EQ&=1-a_c'(C-y)(-y-Cv),
 \notag\\
 MZ-ZM&=a_c-a_c'y\frac{k}{x},
 &QZ-ZQ&=1-a_c'y\frac{h}{x}.
 \label{eq:opt:ordinary-commutators}
\end{align}
They are nonnegative.  Indeed, in the varying region
$h/x=y+Cv\in[y,0]$ and $k/x\in[0,1]$.  With $t=-y\in[0,1]$,
\begin{equation}
 a_c'(C-y)t=\frac{t(C+t)}{c+2t}\le1,
 \qquad
 a_c't^2=\frac{t^2}{c+2t}\le1;
 \label{eq:opt:ordinary-commutator-signs}
\end{equation}
the first inequality is equivalent to
$(t-1)(t+c)\le0$.  In the flat region $a_c'=0$.  Hence $Z,E$ can be
moved to the tail, and a worst word has the form
\begin{equation}
 U^\mu\{M,Q\}^*\{Z,E\}^*.
 \label{eq:opt:canonical-word}
\end{equation}

On the middle block use the simpler potential
\[
 \widehat W(x,N,D)=xD+x^2f_c\!\left(\frac{D-CN}{x}\right).
\]
The identity
\begin{equation}
 2f_c-yf_c'
 =\max\{0,-Cf_c',-Cf_c'+a_c(1+y),
              1-cf_c'+a_cy\}
 \label{eq:opt:ordinary-HJB}
\end{equation}
is a direct rewriting of \eqref{eq:opt:f-hjb-identities}; it pays the
entire ordinary continuation.  If the cap prefix has mass $\mu$, then
after it $x=1-\mu$, $N=D=\mu$, and
$y=-c\mu/(1-\mu)$.  Its internal cap-pair reward plus its worst
continuation is therefore
\begin{equation}
 F_{c,u}(\mu)
 =\frac{2-c(u-1)}2\mu^2+\mu(1-\mu)
 +(1-\mu)^2f_c\!\left(-\frac{c\mu}{1-\mu}\right).
 \label{eq:opt:prefix-objective}
\end{equation}
For $0<\mu<1/C$, differentiation gives
\begin{equation}
 F_{c,u}'(\mu)
 =c\mu\,[Z(c)-u-\operatorname{artanh}\mu].
 \label{eq:opt:prefix-derivative}
\end{equation}
On the mixed curve the unique maximizer in this range is therefore
$\mu=m$, and
\begin{equation}
 F_{c,u}(m)
 =f_c(0)+\frac c2(\operatorname{artanh}m-m)=\frac c2.
 \label{eq:opt:prefix-maximum}
\end{equation}
Here $m\le1/C$, since $c,m\le1/\varphi$.  On the flat part
$\mu\ge1/C$,
\begin{equation}
 F_{c,u}(\mu)=\mu-\frac{c(u-1)}2\mu^2.
 \label{eq:opt:prefix-flat}
\end{equation}
Its derivative at $1/C$ is
\[
 \frac{c(1+m)-1}{C}\le0,
\]
and only decreases thereafter.  Thus $F_{c,u}(\mu)\le c/2$ on the whole
unit interval.  At $c=r,m=0,u=z_*$, the same calculation is maximized at
$\mu=0$, and increasing $u$ only lowers
\eqref{eq:opt:prefix-objective}.

To make this independent exchange proof finite, fix
$\varepsilon>0$ and first replace the last
$\varepsilon n+O(1)$ outcomes by genuine zeros.  Earlier states and
actions are unchanged, and only $O(\varepsilon n^2+n)$ pair terms are
affected.  All cap and ordinary exchanges in the retained prefix then
occur with $x\ge\varepsilon$; their displayed fluid commutators have
uniform Taylor remainders there.  More precisely, after assigning mass
$1/n$ to every job, a normalized adjacent exchange has remainder
$O_{u,\varepsilon}(n^{-3})$; there are at most $n^2$ exchanges.  At the
join $y=-1$, split an exchange at the join and use the two one-sided
formulae, whose values agree.  The artificial zeros already form the
terminal block, so no unbubbled suffix cap invalidates
\eqref{eq:opt:ordinary-commutators}.  The normalized one-job remainder in
the continuation-potential calculation is
$O_{u,\varepsilon}(n^{-2})$.  For fixed $\varepsilon$, first let
$n\to\infty$, and then let $\varepsilon\downarrow0$.  Equations
\eqref{eq:opt:ordinary-HJB}--\eqref{eq:opt:prefix-maximum} give
\[
 \Gamma_n\le c\binom n2+o(n^2),
\]
which is an independent size-asymptotic certificate for the two upper
bounds above.
}
